\section{Standard errors and normal approximations}

A statistic varies from sample to sample. The standard deviation of its sampling
distribution measures the typical size of this variation.

\begin{definitionbox}
The \textbf{standard error} of a statistic \(T\) is the standard deviation of its
sampling distribution:
\[
\operatorname{SE}(T)=\sqrt{\Var(T)}.
\]
\end{definitionbox}

The term does not describe a computational mistake. It describes variation
created by random sampling.

\subsection*{Standard deviation and standard error}

The population standard deviation describes variation among individual
observations. The standard error describes variation among statistics computed
from repeated samples.

If individual observations have standard deviation \(\sigma\), then
\[
\operatorname{SD}(X_i)=\sigma,
\qquad
\operatorname{SE}(\overline X)=\frac{\sigma}{\sqrt n}.
\]
A population may remain highly variable while its mean is estimated much more
precisely.

\subsection*{Standard errors of common statistics}

For independent observations with common variance \(\sigma^2\),
\[
\operatorname{SE}(\overline X)=\frac{\sigma}{\sqrt n}.
\]
For a sample proportion based on independent Bernoulli observations,
\[
\operatorname{SE}(\widehat p)
=
\sqrt{\frac{p(1-p)}{n}}.
\]
Both formulas contain the scale \(1/\sqrt n\). Multiplying the sample size by four
divides the standard error by two.

\begin{remarkbox}
Larger samples reduce random sampling variation, but they do not automatically
remove bias, dependence, measurement error, or model misspecification.
\end{remarkbox}

\subsection*{Normal approximation in standard-error units}

The central limit theorem suggests that, for sufficiently large samples,
\[
\frac{\overline X-\mu}{\operatorname{SE}(\overline X)}
\approx N(0,1).
\]
Equivalently,
\[
\overline X
\approx
\mu+\operatorname{SE}(\overline X)Z,
\qquad Z\sim N(0,1).
\]
The same form applies to a sample proportion:
\[
\frac{\widehat p-p}{\operatorname{SE}(\widehat p)}
\approx N(0,1).
\]
A numerical difference should therefore be interpreted relative to its standard
error.

\subsection*{Typical normal ranges}

For \(Z\sim N(0,1)\),
\[
P(-1.96\leq Z\leq1.96)\approx0.95.
\]
Consequently, when the approximation is appropriate,
\[
P\left(
\theta-1.96\operatorname{SE}(T)
\leq T\leq
\theta+1.96\operatorname{SE}(T)
\right)
\approx0.95,
\]
provided that the sampling distribution of \(T\) is centred at \(\theta\).
This is a statement about the random statistic before the sample is observed.

\subsection*{Known and estimated standard errors}

The theoretical standard error may contain unknown parameters. For the mean,
\[
\operatorname{SE}(\overline X)=\frac{\sigma}{\sqrt n}.
\]
The population standard deviation is commonly estimated using
\[
S^2=\frac1{n-1}\sum_{i=1}^n(X_i-\overline X)^2,
\]
so that
\[
\widehat{\operatorname{SE}}(\overline X)=\frac{S}{\sqrt n}.
\]
The denominator \(n-1\) is used because, under independent identical sampling
with finite variance,
\[
\E[S^2]=\sigma^2.
\]

For a proportion, the plug-in estimate is
\[
\widehat{\operatorname{SE}}(\widehat p)
=
\sqrt{\frac{\widehat p(1-\widehat p)}{n}}.
\]
An estimated standard error is itself random. It varies from sample to sample and
approximates the theoretical standard error only under the stated assumptions.

\subsection*{From standardized distances to inference}

For a proposed parameter value \(\theta_0\),
\[
Z_{\mathrm{obs}}
=
\frac{T_{\mathrm{obs}}-\theta_0}
{\operatorname{SE}_{\theta_0}(T)}
\]
measures the observed difference in standard-error units. Values near zero are
ordinary under the proposed model; values far into a tail are unusual.

A standardized sampling statement may also be inverted to form a range of
parameter values, provided that the dependence of the standard error on the
parameter is handled correctly. This leads to confidence intervals.

\begin{summarybox}
A standard error is the standard deviation of a sampling distribution. It is the
natural scale for comparing an observed statistic with a parameter value, and it
connects limit theorems with confidence intervals and hypothesis tests.
\end{summarybox}
